$h|_{T^*M} = G$ identically (Definition \claimref{obj.sheaf.collar-data-coupling}), so the inclusion $T^*M \hookrightarrow T \subset \shC$ satisfies the Clifford relation of $G$ and extends, by the universal property fibrewise (\claimref{std-algebra.existence-uniqueness}), to an algebra map $\varphi : \shA \to \shC$, smooth in blade frames since it sends blades to blades. Injectivity is Lemma \claimref{sheaf.blade-basis-degenerate-forms-included} applied twice: extend a local frame $(e_1, \dots, e_n)$ of $T^*M$ by $(\nO, \nInf)$ to a frame of $T$; the $2^{n+2}$ blades of the extended frame are a basis of $\shC$'s fibre, the images under $\varphi$ of $\shA$'s blades are exactly the $2^n$ blades in the first $n$ letters, a subset of that basis, hence linearly independent. So $\varphi$ is a global isomorphism onto the subalgebra generated by $T^*M$, over every point, degenerate or not.

Over $\Mo \cap \Mp$ the form $h$ is nondegenerate on $T^*M$ (it restricts to $G$), so $T = T^*M \perp (T^*M)^{\perp_h}$ with the complement of rank two; comparing determinants via Lemma \claimref{sheaf.determinant-identity-blow-up-expansion}(1), $\det h / \det G = -(1 + 2\mu\nu\beta^2 s) < 0$ on the physical collar side (where the factor is $\geq 1$ in absolute value and negative) and $= -1$ outside the collar, so the complement plane carries signature $(1,1)$ (\claimref{std-algebra.sylvester-inertia}): a hyperbolic plane. The embedding $\Cl(p,q) \cong A(x) \subset C(x) \cong \Cl(p{+}1, q{+}1)$ (the fibre of \claimref{sheaf.one-sided-completion}) is therefore exactly the Rule-1 conformal embedding of \cite[\S4.3]{ButlerT1}, performed on the coupled plane in place of the naive one; outside the collar it is literally $G \perp H$.

Over $x \in \Sig$, set $k^* := \nO/(2\nu)$. Then $(k^*)^2 = \nO^2/(4\nu^2) = 0$ and $k k^* + k^* k = 2\,h(k, k^*) = 2 \cdot \nu/(2\nu) = 1$ by Lemma \claimref{sheaf.wall-gram}(2): a ladder pair. Off $\Sig$ on the physical side, $k^2 = G(k,k) = \lambda \neq 0$, so $k$ is not even null there and no ladder relation can hold; the ideal-theoretic contrast on $\Sig$ is Theorem \claimref{sheaf.degeneration-sheaf} against Proposition \claimref{sheaf.resolution-ideal-theoretically}, both quoted in the statement. The pair $(k, k^*)$ satisfies the ladder-pair relations of \cite[\S7]{ButlerT1} (the canonical anticommutation relations of one fermionic mode, in the naming of the tower paper \cite{ButlerT3}) exactly on $\Sig$, which is the localisation claim.
